% !TEX root = ../main.tex
\chapter{AURORA: o visualizador de RTL PRISM}
\label{cap:11-aurora-prism-rtl}

O \tool{PRISM} é o visualizador de RTL (\emph{register-transfer level}) embarcado
na \tool{AURORA}. A partir do Verilog sintetizável de um projeto SAPHO, ele
produz um \emph{esquemático} navegável do circuito: cada módulo da hierarquia
vira um diagrama vetorial (SVG) com células, portas e fios, no qual o usuário
pode entrar em submódulos com um clique, voltar ao código-fonte com duplo-clique
e --- desde a fase O9 --- alternar para uma \emph{simulação lógica interativa}
do mesmo \emph{netlist} via \tool{DigitalJS}.

Este capítulo documenta o subsistema completo: a janela e o canal de IPC
(\file{main/ipc/prism.js}), o \emph{preload} sandboxizado
(\file{js/app/preload\_prism.js}), o renderer da janela
(\file{html/prism/prism.js} e \file{html/prism/prism.html}), o \emph{pipeline}
Verilog~$\rightarrow$~Yosys~$\rightarrow$~netlistsvg-aurora~$\rightarrow$~SVG, o
fork \repo{netlistsvg-aurora} (o que ele muda em relação ao
\lstinline|@silimate/netlistsvg| de \emph{upstream}), o sistema de \emph{skins}
customizadas da família SAPHO, e o estado de integração de
\tool{DigitalJS}~+~\tool{yosys2digitaljs}. A síntese pela ferramenta \tool{Yosys}
e o \emph{toolchain} bundle são tratados em outros capítulos; aqui o foco é o que
o \tool{PRISM} faz com a saída do \tool{Yosys}.

\begin{nota}
Fonte da verdade: todas as afirmações deste capítulo foram verificadas lendo o
código real. Onde a documentação de prosa interna e o código divergem, vale o
código. Os arquivos efetivamente abertos estão listados ao final do estudo.
\end{nota}

\section{Visão geral e fluxo de acionamento}

O \tool{PRISM} é acionado pelo botão \enquote{PRISM} da \emph{toolbar} de
compilação da \tool{AURORA} (id \lstinline|prismcomp|). O \emph{handler} desse
botão vive em \file{js/compilation/compilation\_flow.js}, na função
\lstinline|handlePrismStep()|. Conforme a filosofia \enquote{cada botão é
self-contained} desse arquivo, o \tool{PRISM} é um \emph{superset} do botão
Verilog: ele primeiro roda \cmm{}~+~ASM para cada processador conhecido e a
verificação de sintaxe Verilog (top-level), e só então invoca o \tool{Yosys}
para a análise estrutural.

\begin{lstlisting}[language=bash,caption={Expansão dos botões (cabeçalho de compilation\_flow.js)}]
Verilog = cmm + asm + iverilog(-tnull, top)
PRISM   = Verilog                            + yosys
Wave    =     cmm + asm + iverilog(-o vvp, tb) + vvp + gtkwave
\end{lstlisting}

O botão \lstinline|prismcomp| só fica habilitado quando o \file{.spf} tem um
\emph{top-level} definido (\lstinline|setEnabled('prismcomp', hasTop)| em
\file{compilation\_flow.js}). O fluxo de alto nível é:

\begin{enumerate}[leftmargin=2.2em]
  \item \textbf{Renderer (Aurora principal):} \lstinline|handlePrismStep()| roda
        \cmm{}/ASM e \lstinline|verilogSyntaxCheck()|, monta o \emph{payload} de
        caminhos absolutos (\lstinline|buildPrismCompilationPaths|), consulta o
        \emph{store} de \emph{overrides} de IA para a etapa
        \lstinline|prism-yosys| e chama \lstinline|electronAPI.prismCompileWithPaths(paths)|.
  \item \textbf{Main:} o \emph{handler} IPC \lstinline|prism-compile-with-paths|
        executa o \emph{pipeline} de síntese e geração de SVG, e abre (ou foca) a
        janela do \tool{PRISM}.
  \item \textbf{Renderer (janela PRISM):} recebe \lstinline|compilation-complete|
        com o caminho do SVG do top-level, carrega o SVG, instrumenta cliques e
        navegação.
\end{enumerate}

O \emph{payload} de caminhos é montado em
\lstinline|buildPrismCompilationPaths(projectPath)| com todos os campos
absolutos e \emph{slashes} normalizadas para barra direta:

\begin{lstlisting}[caption={Campos do payload de caminhos (compilation\_flow.js)}]
{
  projectPath, componentsPath,
  hdlPath:      <components>/HDL,
  tempPath:     <components>/Temp/PRISM,
  yosysPath:    <components>/Packages/msys/mingw64/bin/yosys.exe,
  spfPath:      <currentSpfPath>,
  topLevelPath: <projectPath>/TopLevel,
  yosysOverride?: { appendArgs, prependArgs, removeArgs, envSet, envUnset }
}
\end{lstlisting}

O \emph{main} também expõe \lstinline|get-prism-compilation-paths|, que recompõe
esses mesmos campos a partir do \file{.spf} aberto
(\lstinline|state.currentOpenProjectPath|); a janela usa esse canal para
recompilar e para construir a simulação \tool{DigitalJS} por conta própria.

\section{A janela do PRISM}

A janela é criada por \lstinline|createPrismWindow(compilationData)| em
\file{main/ipc/prism.js}. É uma \lstinline|BrowserWindow| única
(\lstinline|state.prismWindow|): se já existe e não foi destruída, a chamada
apenas a foca e reenvia o payload de compilação, em vez de abrir uma segunda
janela.

\subsection{Configuração da BrowserWindow}

\begin{tabularx}{\linewidth}{Y Z}
\toprule
\textbf{Propriedade} & \textbf{Valor / efeito} \\
\midrule
\lstinline|width| / \lstinline|height| & 1400 $\times$ 900 (mín. 1000 $\times$ 700) \\
\lstinline|frame| & \lstinline|false| --- janela \emph{frameless} com \emph{titlebar} customizada desenhada pelo HTML do PRISM \\
\lstinline|thickFrame| & \lstinline|true| --- preserva \emph{snap} do Aero e redimensionamento por borda nativo no Windows \\
\lstinline|titleBarStyle| & \lstinline|'hidden'| \\
\lstinline|backgroundColor| & \lstinline|'#0A0D14'| (mesmo fundo escuro da \tool{AURORA}) \\
\lstinline|icon| & \path{assets/icons/sapho_aurora_icon.ico} \\
\lstinline|contextIsolation| & \lstinline|true| \\
\lstinline|sandbox| & \lstinline|true| \\
\lstinline|nodeIntegration| & \lstinline|false| \\
\lstinline|preload| & \path{js/app/preload_prism.js} (existência verificada antes de criar a janela) \\
\bottomrule
\end{tabularx}

A página em si é carregada por \lstinline|loadPage(prismWindow, 'html/prism/prism.html')|,
de \file{main/render\_loader.js}. Esse \emph{loader} é compartilhado por todas as
janelas de primeira parte e resolve a fonte do renderer em duas vias: servidor de
desenvolvimento Vite (quando \lstinline|AURORA_RENDERER_URL| está setado e o app
não está empacotado) ou o \emph{bundle} de produção em \path{dist/} via
\lstinline|file://|. \emph{Não há \emph{fallback}} para o fonte cru --- um
\emph{bundle} ausente é tratado como erro de \emph{build} (página de erro
explícita), porque pós-migração para \tool{Lit} o HTML cru carrega \emph{imports}
que só resolvem através do empacotador.

\subsection{Controles de janela e estado}

Como a janela é \emph{frameless}, os botões minimizar/maximizar/fechar são
desenhados em SVG no HTML e ligados via IPC genéricos do \emph{preload}
(\lstinline|window:minimize|, \lstinline|window:maximize-toggle|,
\lstinline|window:close|). O \emph{main} relê o estado de maximização/tela-cheia
para o renderer pelo canal \lstinline|prism:window-state| nos eventos
\lstinline|maximize|, \lstinline|unmaximize|, \lstinline|enter-full-screen| e
\lstinline|leave-full-screen|, para que o ícone \emph{maximize/restore} se
atualize. Ao carregar, a janela é maximizada e exibida; o \emph{main} também
notifica a janela principal com \lstinline|prism-status| (\lstinline|true| ao
abrir, \lstinline|false| ao fechar).

Erros de carregamento são tratados em três pontos: o \lstinline|try/catch| em
torno de \lstinline|loadPage| (mostra \lstinline|dialog.showMessageBox| e destrói
a janela), o \lstinline|did-fail-load| e o \lstinline|render-process-gone|
(ambos registram em \lstinline|electron-log| e exibem diálogo de erro).

\section{O preload sandboxizado}

\file{js/app/preload\_prism.js} expõe um \emph{único} objeto
\lstinline|window.electronAPI| via \lstinline|contextBridge|, com uma
\emph{allowlist enumerada} de canais. O comentário no próprio arquivo é
explícito: \emph{deliberadamente não há} um \lstinline|send|/\lstinline|removeAllListeners|
genérico, para que esse renderer nunca alcance um canal IPC arbitrário.

\begin{tabularx}{\linewidth}{Y Z}
\toprule
\textbf{Método exposto} & \textbf{Canal IPC / efeito} \\
\midrule
\lstinline|windowMinimize/MaximizeToggle/Close| & \lstinline|window:minimize| / \lstinline|maximize-toggle| / \lstinline|close| \\
\lstinline|onWindowState(cb)| & escuta \lstinline|prism:window-state| \\
\lstinline|readFile(path)| & \lstinline|read-file| (passa por \lstinline|safePath()| no main; evita \lstinline|fetch('file://')|) \\
\lstinline|joinPath(...segs)| & \lstinline|join-path| \\
\lstinline|getPrismCompilationPaths()| & \lstinline|get-prism-compilation-paths| \\
\lstinline|prismCompileWithPaths(paths)| & \lstinline|prism-compile-with-paths| \\
\lstinline|prismRecompile(paths)| & \lstinline|prism-recompile| \\
\lstinline|buildDigitalJS(paths, module)| & \lstinline|prism:build-digitaljs| (O9) \\
\lstinline|exportWave(payload)| & \lstinline|prism:export-wave| (grava o \lstinline|.vcd| do monitor no Temp do PRISM e manda a janela principal abri-lo, via \lstinline|aurora:open-wave|) \\
\lstinline|generateSVGFromModule(name, dir)| & \lstinline|generate-svg-from-module| \\
\lstinline|openSourceFile({filePath,line,column})| & \lstinline|prism:open-source-file| \\
\lstinline|onCompilationComplete(cb)| & escuta \lstinline|compilation-complete| \\
\lstinline|onTerminalLog| / \lstinline|onPrismStatus| & escuta \lstinline|terminal-log| / \lstinline|prism-status| \\
\bottomrule
\end{tabularx}

Além disso, o \emph{preload} reencaminha eventos \lstinline|terminal-log| do
\emph{main} para a página como \lstinline|window.postMessage|, de modo que a
janela do \tool{PRISM} possa exibir o progresso da síntese.

\section{O pipeline de compilação}

O coração do \tool{PRISM} é a função assíncrona
\lstinline|performPrismCompilationWithPaths(compilationPaths)|, que orquestra
três etapas: síntese com \tool{Yosys}, divisão do \emph{netlist} por módulo, e
geração do SVG do top-level. O top-level vem \emph{sempre} de
\lstinline|spf.structure.topLevelFile| (\lstinline|path.basename(...,'.v')|).

\begin{lstlisting}[caption={Diagrama de fluxo do pipeline (main/ipc/prism.js)}]
 +-------------------------+
 |  spf.structure.topLevel |  ->  topLevelModule (basename sem .v)
 +-----------+-------------+
             |
             v
 +-------------------------+    coleta unificada de .v + dedup por path
 |  runYosysCompilation    |    -> yosys_script.ys -> yosys.exe -> hierarchy.json
 +-----------+-------------+
             |
             v
 +-------------------------+    para cada modulo "clicavel":
 |  splitHierarchyJson     |    limpa nome, escreve <modulo>.json no tempDir
 +-----------+-------------+
             |
             v
 +-------------------------+    skin default + skins custom -> netlistsvg.render
 |  generateModuleSVG      |    -> injeta data-cell-type -> <modulo>.svg
 +-----------+-------------+
             |
             v
 +-------------------------+
 |  janela PRISM carrega   |  ->  compilation-complete { topLevelModule, svgPath, tempDir }
 +-------------------------+
\end{lstlisting}

\subsection{Etapa 1 --- coleta de arquivos e síntese com Yosys}

A função \lstinline|runYosysCompilationWithPaths(paths, topLevelModule, tempDir)|
coleta \emph{todos} os \file{.v} relevantes em um \lstinline|Set| (dedup por
caminho absoluto, evitando erros de módulo duplicado no \tool{Yosys}). As fontes,
nesta ordem, são:

\begin{enumerate}[leftmargin=2.2em]
  \item \path{components/HDL/*.v} --- a biblioteca SAPHO (processor, addr\_dec,
        core, ula, myFIFO, instr\_dec etc.), sempre incluída porque qualquer
        \emph{top} que seja um processador SAPHO depende dela. Arquivos com
        \lstinline|_tb| ou \lstinline|test| no nome são excluídos.
  \item \lstinline|spf.structure.synthesizableFiles[].path| --- a fonte canônica
        de todos os \file{.v} do projeto (auto-descobertos pela árvore de
        arquivos e/ou adicionados manualmente). Caminhos relativos são resolvidos
        contra \lstinline|spf.structure.basePath| (formato novo do \file{.spf};
        caminhos absolutos passam direto).
  \item \path{<projeto>/TopLevel/*.v}, se a pasta existir (\emph{wrapper}
        top-level, raro mas possível).
  \item \textbf{\emph{Fallback} defensivo:} para cada processador listado em
        \lstinline|spf.structure.processors|, varre
        \path{<projeto>/<proc>/Hardware/*.v}. No caminho feliz o \emph{dedup}
        torna isso \emph{no-op}, mas garante os módulos do processador em
        projetos antigos onde \lstinline|synthesizableFiles| esteja vazio.
\end{enumerate}

Se nenhum \file{.v} for encontrado, a etapa falha com \enquote{No Verilog files
found for compilation}.

O \emph{script} \tool{Yosys} é gerado em \path{<tempDir>/yosys_script.ys}. Cada
arquivo é lido com \lstinline|read_verilog -setattr src| --- a opção
\lstinline|-setattr src| faz o \tool{Yosys} gravar um atributo
\lstinline|src="arquivo.v:linha.col-linha.col"| em cada célula derivada do
\emph{source}, o que mais tarde habilita o duplo-clique \enquote{abrir no
editor}. O \emph{script} é exatamente:

\begin{lstlisting}[caption={Script Yosys do esquemático (gerado em main/ipc/prism.js)}]
read_verilog -setattr src "arquivo1.v"
read_verilog -setattr src "arquivo2.v"
...
hierarchy -top <topLevelModule>
proc
setundef -zero
opt_clean -purge
write_json "hierarchy.json"
\end{lstlisting}

\begin{tabularx}{\linewidth}{Y Z}
\toprule
\textbf{Comando} & \textbf{Função no PRISM} \\
\midrule
\lstinline|read_verilog -setattr src| & lê cada \file{.v} e anota a procedência (\lstinline|src|) em cada célula \\
\lstinline|hierarchy -top| & fixa o módulo top e resolve a árvore de instâncias/parametrizações \\
\lstinline|proc| & converte processos (\lstinline|always|) em células (mux, dff etc.) \\
\lstinline|setundef -zero| & troca \emph{don't-care} (\lstinline|x|) por constante 0, evitando \enquote{linhas fantasma} que o netlistsvg desenharia a partir de ramos \emph{full\_case} inalcançáveis \\
\lstinline|opt_clean -purge| & remove fios/células desconectados após a substituição \\
\lstinline|write_json| & emite o \emph{netlist} consolidado em \path{hierarchy.json} \\
\bottomrule
\end{tabularx}

Note que o \emph{script do esquemático não faz} \lstinline|techmap|/\lstinline|abc|:
o objetivo é um diagrama em nível de palavra (RTL), não \emph{gates}.

O \tool{Yosys} é executado com \lstinline|spawnTracked(yosysExe, finalArgs, ...)|.
O uso de \lstinline|spawnTracked| (de \file{main/process\_registry.js}) garante
que fechar a janela principal mata uma síntese do \tool{PRISM} em andamento ---
necessário porque o \tool{Yosys} roda do \path{mingw64/bin} bundle, fora de
\path{Temp/}. Os argumentos base são \lstinline|['-s', yosysScriptPath]|, mas
podem ser modificados por um \emph{override} de IA
(\lstinline|paths.yosysOverride|) com o mesmo contrato dos demais \emph{steps}:
\lstinline|removeArgs|, \lstinline|prependArgs|, \lstinline|appendArgs|,
\lstinline|envSet|, \lstinline|envUnset|. Sem \emph{override}, o comportamento é
idêntico ao padrão. O processo só resolve em sucesso se o código de saída for 0
\emph{e} \path{hierarchy.json} existir; caso contrário rejeita com o
\lstinline|stderr| capturado.

\subsection{Etapa 2 --- divisão do netlist por módulo}

\lstinline|splitHierarchyJson(hierarchyJsonPath, tempDir)| lê o
\path{hierarchy.json} consolidado e, para cada módulo \enquote{clicável}, escreve
um JSON \emph{por módulo} (\path{<sanitized>.json}) contendo apenas aquele
módulo. Isso é o que permite, em tempo de navegação, renderizar qualquer módulo
sob demanda sem reexecutar o \tool{Yosys}.

Dois \emph{helpers} regem o que entra:

\begin{description}[leftmargin=1.4em,style=nextline]
  \item[\lstinline|isClickableModule(name)|] descarta primitivas do \tool{Yosys}
        (qualquer \lstinline|$_|, \lstinline|$dff|, \lstinline|$mux|,
        \lstinline|$add|, \lstinline|$pmux|, \lstinline|$lut|, etc.) e aceita
        apenas \lstinline|$paramod...| ou identificadores Verilog legítimos
        (\lstinline|^[a-zA-Z_][a-zA-Z0-9_]*$|). É o mesmo conjunto de padrões
        replicado no renderer (\lstinline|_isClickableModuleType|), para que o
        que o \emph{main} grava como navegável coincida com o que o renderer
        torna clicável.
  \item[\lstinline|cleanModuleName(name)|] normaliza nomes gerados pelo
        \tool{Yosys}: desembrulha prefixos \lstinline|$paramod\...| de módulos
        parametrizados, remove \emph{hashes} (\lstinline|$<hex de 40+ dígitos>|),
        sufixos de parâmetro (\lstinline|\PARAM=...|), e prefixos de escopo de
        \emph{generate block} (\lstinline|genblk<N>.|), de modo que o usuário veja
        o nome real (ex.: \lstinline|my_f2i| em vez de
        \lstinline|genblk27.my_f2i|).
\end{description}

Para cada módulo aceito, o nome é limpo e \emph{sanitizado}
(\lstinline|sanitizeFileName|), as células internas também têm seus nomes e tipos
limpos, e o JSON resultante é escrito com a forma
\lstinline|{ creator, modules: { <cleanName>: <data> } }|.

\subsection{Etapa 3 --- renderização do SVG com netlistsvg-aurora}

\lstinline|generateModuleSVGWithPaths(moduleName, tempDir)| produz o SVG de um
módulo. Pontos-chave:

\begin{itemize}[leftmargin=1.4em]
  \item A biblioteca é carregada \emph{in-process} via
        \lstinline|require('@silimate/netlistsvg')| --- \emph{não há} \emph{spawn}
        de processo externo nem binário separado para o netlistsvg. (O comentário
        em \lstinline|get-prism-compilation-paths| confirma: \enquote{netlistsvg
        agora vem do node\_modules e roda in-process}.)
  \item A \emph{skin} é obtida por \lstinline|getDefaultSkinData()| (ver
        \autoref{sec:skins}), que combina a \path{lib/default.svg} do pacote com
        as \emph{skins} customizadas de \path{assets/prism-skins/}.
  \item Antes de renderizar, \lstinline|pruneNetlistToSkinPorts| poda do
        \emph{netlist} conexões a portas que a \emph{skin} customizada do tipo
        \emph{não} desenha (ver \autoref{sec:skin-ports}), evitando que o ELK
        aborte o \emph{layout}.
  \item A renderização usa \lstinline|netlistsvgLib.render(skinData, netlistJson, cb)|
        (estilo \emph{callback}), envolto em \lstinline|Promise|.
  \item Se a renderização não devolver uma \emph{string} não-vazia, a função
        falha com mensagem clara (em vez de estourar \lstinline|TypeError| na
        injeção seguinte).
  \item Por fim, \lstinline|injectCellTypesIntoSvg| insere
        \lstinline|data-cell-type=<tipo>| em cada \lstinline|<g id="cell_<inst>">|,
        e o SVG é gravado em \path{<sanitized>.svg}.
\end{itemize}

A injeção de \lstinline|data-cell-type| é necessária porque o netlistsvg rotula
cada célula com o \emph{nome da instância} (\lstinline|s:attribute="ref"|), não
com o \emph{tipo} de módulo. \lstinline|buildInstanceTypeMap| constrói o mapa
\lstinline|instância -> tipo| a partir do JSON do \tool{Yosys}, e
\lstinline|injectCellTypesIntoSvg| usa uma \emph{regex} (não um \emph{parser} XML)
para anotar cada \lstinline|<g>| de célula com o tipo --- assim o renderer sabe
para qual módulo navegar ao clicar. Valores de atributo são escapados por
\lstinline|xmlAttrEscape| (\lstinline|&|, \lstinline|"|, \lstinline|<|).

\subsection{Resultado e abertura da janela}

Em sucesso, \lstinline|performPrismCompilationWithPaths| retorna
\lstinline|{ success: true, topLevelModule, svgPath, tempDir }|, e o \emph{handler}
\lstinline|prism-compile-with-paths| chama \lstinline|createPrismWindow(result)|.
A janela recebe \lstinline|compilation-complete| com esse objeto, carrega o SVG
do top-level e inicializa a navegação. Durante todo o processo, mensagens de
progresso são enviadas à janela principal pelo canal \lstinline|terminal-log| no
terminal \lstinline|tveri| (\enquote{Top-level: \dots}, \enquote{Running Yosys
synthesis\dots}, \enquote{PRISM compilation completed successfully}, ou os erros
correspondentes).

\section{O fork netlistsvg-aurora}
\label{sec:fork}

O \file{package.json} da \tool{AURORA} aponta a dependência
\lstinline|@silimate/netlistsvg| para o fork do laboratório:

\begin{lstlisting}[caption={Dependência apontando para o fork (package.json da AURORA)}]
"@silimate/netlistsvg":
  "github:nipscernlab/netlistsvg-aurora#58a0703904a35841616db58716f40532b84fb018"
\end{lstlisting}

Ou seja, embora o nome do pacote permaneça \lstinline|@silimate/netlistsvg|
(versão \lstinline|1.1.4|, autor original Neil Turley), o código instalado vem de
\repo{netlistsvg-aurora}, fixado por \emph{commit hash}.

\subsection{O que o fork muda em relação ao upstream}

A diferença substantiva do fork é um \textbf{ajuste de \emph{tsconfig} para
destravar a compilação sob TypeScript~6}, seguida de uma reconstrução de
\path{built/}. O \file{tsconfig.json} do fork é:

\begin{lstlisting}[caption={tsconfig.json do fork netlistsvg-aurora}]
{
  "compilerOptions": {
    "outDir": "./built",
    "allowJs": true,
    "lib": [ "es6", "dom" ],
    "target": "es2015",
    "module": "commonjs",
    "moduleResolution": "node",
    "types": []
  },
  "include": [ "./lib/*", "./lib/@types" ]
}
\end{lstlisting}

A mensagem de \emph{commit} do fork descreve precisamente as mudanças:

\begin{itemize}[leftmargin=1.4em]
  \item \lstinline|target| passou de \lstinline|"es5"| (que o TS~6 marca como
        \emph{deprecated} e recusa compilar, erro \lstinline|TS5107|) para
        \lstinline|"es2015"|.
  \item Foram adicionados \lstinline|module: "commonjs"| e
        \lstinline|moduleResolution: "node"| explícitos, para que a sintaxe
        CommonJS existente (\lstinline|import x = require(...)|) continue
        funcionando sob o compilador moderno.
  \item \lstinline|types: []| foi fixado para que o \lstinline|tsc| não varra
        pacotes \lstinline|@types| transitivos e tropece em definições de
        \emph{babel}/\emph{jest}.
  \item O diretório \path{built/} foi reconstruído a partir dessas
        configurações; a mensagem afirma que a saída é funcionalmente idêntica
        (testada renderizando os exemplos digitais).
\end{itemize}

O \emph{layout} de pastas do fork (\path{lib/} com os \file{.ts} fonte ---
\file{index.ts}, \file{FlatModule.ts}, \file{Cell.ts}, \file{Port.ts},
\file{Skin.ts}, \file{drawModule.ts}, \file{elkGraph.ts}, \file{YosysModel.ts}
--- além de \file{lib/default.svg} e \file{lib/analog.svg}, e \path{built/} com os
\file{.js} compilados e \file{netlistsvg.bundle.js}) é o mesmo do \emph{upstream}.
A API usada pela \tool{AURORA} é a função
\lstinline|render(skinData, yosysNetlist, done)| de \file{lib/index.ts}, que
constrói o \lstinline|FlatModule|, monta o grafo ELK
(\lstinline|buildElkGraph|), roda o \emph{layout} via \lstinline|elkjs| e desenha
com \lstinline|drawModule| --- preservando o estilo de \emph{callback} legado que
o código do \tool{PRISM} consome.

\begin{nota}
As demais diferenças efetivas em relação ao netlistsvg de prateleira estão
\emph{fora} do fork da biblioteca: a \tool{AURORA} não edita
\path{lib/default.svg}; em vez disso \emph{mescla} suas próprias \emph{skins} em
tempo de execução (ver \autoref{sec:skins}). As \emph{skins} customizadas da
família SAPHO vivem no repositório \tool{AURORA}, em \path{assets/prism-skins/}.
\end{nota}

\section{Sistema de skins}
\label{sec:skins}

Uma \emph{skin} do netlistsvg é um SVG que declara, em blocos
\lstinline|<g s:type="...">|, como desenhar cada \emph{tipo} de célula (suas
portas, identificadas por \lstinline|s:pid|, e a geometria do símbolo). A
\path{lib/default.svg} do pacote traz os tipos genéricos: portas externas
(\lstinline|inputExt|, \lstinline|outputExt|), \emph{flip-flops} e \emph{latches}
(\lstinline|dff|, \lstinline|dffn|, \lstinline|dlatch| e variantes \lstinline|-bus|),
operadores aritméticos e lógicos (\lstinline|add|, \lstinline|sub|, \lstinline|mul|,
\lstinline|div|, \lstinline|mod|, \lstinline|and|, \lstinline|or|, \lstinline|not|,
\lstinline|mux|, \lstinline|eq|, \lstinline|lt|, \lstinline|shl|, \lstinline|shr|
etc.), \lstinline|split|/\lstinline|join|, \lstinline|constant| e o fallback
\lstinline|generic|.

\subsection{Mescla de skins customizadas em runtime}

\lstinline|getDefaultSkinData()| (em \file{main/ipc/prism.js}) não usa a
\path{default.svg} crua: ela a combina com as \emph{skins} de
\path{assets/prism-skins/}. O processo, executado a \emph{cada} compilação (sem
\emph{cache}, para que \enquote{Recompile} já reflita edições em \emph{skins}
custom):

\begin{enumerate}[leftmargin=2.2em]
  \item Lê a \path{lib/default.svg} resolvida via
        \lstinline|require.resolve('@silimate/netlistsvg')|.
  \item \lstinline|loadCustomSkinBlocks(customSkinDir)| varre
        \path{assets/prism-skins/*.svg} (ignorando arquivos começados por
        \lstinline|_|, como \file{\_template.svg}), e de cada um extrai os blocos
        \lstinline|<g s:type="...">| de topo via
        \lstinline|extractTopLevelGBlocks|. Esse extrator usa contagem de
        profundidade de \lstinline|<g>| (pilha) para suportar \lstinline|<g>|
        aninhados dentro de portas/labels --- uma \emph{regex} \emph{non-greedy}
        quebraria em blocos como o \lstinline|generic|.
  \item Para cada \lstinline|s:type| customizado, o bloco correspondente é
        \emph{removido} da \emph{skin} default e o customizado é injetado logo
        antes de \lstinline|</svg>|.
\end{enumerate}

O comentário no código explica a escolha de mesclar em tempo de execução em vez
de editar \path{node_modules/.../default.svg}: um \lstinline|npm install|
sobrescreveria a \path{default.svg}; mantendo as customizações em
\path{assets/prism-skins/} elas sobrevivem ao \emph{install} e ficam versionadas
no git.

\subsection{A família de símbolos SAPHO}

As \emph{skins} customizadas são geradas por
\file{scripts/prism-skin-standard.js}, que estabelece um \emph{padrão} de
símbolos: cada classe de módulo é desenhada como seu símbolo de livro-texto (mux
= pentágono apontando à direita; ALU = \emph{chevron} entalhado; decodificador =
trapézio \emph{fan-out}; registrador = cartão com \emph{clock} $\blacktriangleright$;
memória = grade; FIFO = \emph{chevrons} de fluxo; blocos hierárquicos/de controle
= cartão de \emph{chip} arredondado). Cada símbolo usa fundo escuro de \emph{chip},
fluxo da esquerda para a direita (dados a oeste, resultados a leste; sinais de
controle como \lstinline|clk|/\lstinline|rst|/\lstinline|op|/\lstinline|en| em
violeta), portas rotuladas no \emph{anchor} e cores temáticas via variáveis CSS
com \emph{fallbacks} em hexadecimal.

Há \emph{skins} para a biblioteca SAPHO inteira --- por exemplo
\file{processor.svg}, \file{core.svg}, \file{ula.svg}, \file{addr\_dec.svg},
\file{instr\_dec.svg}, \file{instr\_fetch.svg}, \file{mem\_ctrl.svg},
\file{mem\_data.svg}, \file{mem\_instr.svg}, \file{myFIFO.svg}, \file{pc.svg},
\file{stack.svg}, \file{prefetch.svg} --- além de uma família grande de variantes
de ALU (\file{ula\_add.svg}, \file{ula\_fadd.svg}, \file{ula\_fmlt.svg},
\file{ula\_i2f.svg}, \file{ula\_mux.svg} etc.). Cada arquivo declara seu
\lstinline|s:type| (o tipo de célula que casa com o nome do módulo Verilog) e um
\lstinline|<s:alias val="..."/>|; uma regra dura do padrão é que cada
\lstinline|s:pid| \emph{deve} ser igual ao nome da porta Verilog.

\begin{lstlisting}[language=Java,caption={Cabeçalho de uma skin SAPHO (assets/prism-skins/processor.svg)}]
<svg xmlns="..." xmlns:s="https://github.com/nturley/netlistsvg">
  <g s:type="processor" transform="translate(0, 0)"
     s:width="150" s:height="132">
    <s:alias val="processor"/>
    ...
    <g s:x="0"   s:y="66" s:pid="clk"/>
    <g s:x="0"   s:y="75" s:pid="rst"/>
    <g s:x="150" s:y="66" s:pid="io_out"/>
    ...
  </g>
</svg>
\end{lstlisting}

\subsection{Poda de portas ausentes na skin}
\label{sec:skin-ports}

Uma \emph{skin} customizada tem um conjunto \emph{fixo} de portas. Se o RTL tem
uma porta que a \emph{skin} não declara (ex.: o \lstinline|core| ganhou um sinal
novo mas \file{core.svg} ainda não tem o \emph{anchor}), o netlistsvg geraria uma
\emph{aresta} para uma \emph{shape} de porta que nunca desenhou, e o ELK abortaria
com \enquote{Referenced shape does not exist}, derrubando o \emph{render} inteiro
do módulo.

Para evitar isso, \lstinline|generateModuleSVGWithPaths| poda o \emph{netlist}
antes de renderizar:

\begin{itemize}[leftmargin=1.4em]
  \item \lstinline|buildSkinPortMap(skinData)| monta
        \lstinline|tipo -> Set(s:pid)| a partir dos blocos da \emph{skin}.
  \item \lstinline|pruneNetlistToSkinPorts(netlistJson, skinPortMap)| remove, de
        cada célula, as conexões cujo \emph{port} não está na \emph{skin} de seu
        tipo (e o respectivo \lstinline|port_directions|), registrando um aviso
        com \lstinline|log.warn| indicando que basta adicionar um
        \lstinline|<g s:pid="..."/>| à \emph{skin} para desenhá-la.
  \item Células de tipo \emph{genérico} (sem \emph{skin}) \emph{não} são tocadas:
        ali o netlistsvg cria as portas a partir das próprias conexões.
\end{itemize}

O resultado é que uma porta extra apenas não aparece (com aviso), em vez de
quebrar a tela.

\section{O renderer da janela}

\file{html/prism/prism.js} (carregado como \lstinline|<script type="module">|)
define a classe \lstinline|PRISMViewer|, instanciada em
\lstinline|DOMContentLoaded| e exposta como \lstinline|window.prismViewer|. Ela
cobre i18n, carregamento e interação do SVG, navegação, pan/zoom e a simulação
\tool{DigitalJS}.

\subsection{Internacionalização}

O renderer mantém um dicionário próprio \lstinline|PRISM_STRINGS| com locais
\lstinline|en| e \lstinline|pt| (não acessa a camada i18n da \tool{AURORA}
principal). \lstinline|detectLocale()| lê \lstinline|aurora-locale| (ou o legado
\lstinline|aurora-yanc-lang|) do \lstinline|localStorage|, com \emph{default}
\lstinline|pt|. As \emph{strings} resolvidas são aplicadas aos elementos da
\emph{toolbar} no carregamento.

\subsection{Carregamento e interação do SVG}

Em \lstinline|compilation-complete|, \lstinline|_onCompilationComplete(data)|
zera as pilhas de navegação, define o módulo corrente e chama
\lstinline|_loadSVG(svgPath, module)|. \lstinline|_loadSVG| lê o SVG via
\lstinline|electronAPI.readFile| (passando pelo \emph{main}, sem
\lstinline|fetch('file://')|), injeta-o no DOM com \emph{fade-in}, atualiza
\emph{breadcrumbs} e info do módulo, e instala as interações.

\lstinline|_setupSVGInteractions()| instala três comportamentos:

\begin{description}[leftmargin=1.4em,style=nextline]
  \item[Navegação por clique] Cada \lstinline|g[data-cell-type]| cujo tipo é
        clicável (mesma lista de \emph{skip} do \emph{main}, em
        \lstinline|_isClickableModuleType|) ganha cursor de ponteiro, realce ao
        passar o mouse e \emph{tooltip}. Um clique simples agenda
        \lstinline|navigateToModule(type)| após $\sim$250\,ms, dando janela para
        um eventual duplo-clique cancelar a navegação.
  \item[Abrir o código-fonte por duplo-clique] Cada
        \lstinline|g[onclick^="gotosrc"]| (atributo que o netlistsvg emite a
        partir do \lstinline|src| do \tool{Yosys}) tem seu \lstinline|onclick|
        removido e substituído por um \emph{listener} de \lstinline|dblclick| que
        extrai o localizador \lstinline|arquivo:linha.col| e chama
        \lstinline|electronAPI.openSourceFile(...)|. (O \emph{stub} global
        \lstinline|window.gotosrc = () => {}| neutraliza o \emph{handler} inline
        original.)
  \item[Realce de fios] Cliques em \lstinline|path|/\lstinline|line|/\lstinline|polyline|
        disparam \lstinline|_highlightWireConnection|, que faz \emph{flood-fill}
        geométrico de \emph{endpoints} (tolerância de 5 unidades) para realçar
        toda a conexão elétrica em \lstinline|var(--aurora-mint)|.
\end{description}

\subsection{Navegação hierárquica}

A entrada em submódulos é feita por \lstinline|navigateToModule(name)|, que chama
\lstinline|electronAPI.generateSVGFromModule(name, tempDir)| --- o \emph{main}
gera o SVG do \path{<modulo>.json} já existente no \emph{tempDir} (sem reexecutar
o \tool{Yosys}). O renderer mantém duas pilhas, \lstinline|navigationHistory| e
\lstinline|forwardHistory|, e oferece:

\begin{itemize}[leftmargin=1.4em]
  \item \lstinline|navigateBack()| / \lstinline|navigateForward()| (também ligados
        aos botões laterais do mouse X1/X2 e a \lstinline|Esc| para voltar);
  \item \emph{breadcrumbs} clicáveis que truncam o histórico até o nível
        escolhido;
  \item \lstinline|recompile()|, que reobtém os caminhos via
        \lstinline|getPrismCompilationPaths| e chama
        \lstinline|prismRecompile|.
\end{itemize}

\subsection{Pan, zoom e download}

O canvas é uma \lstinline|<div>| com \emph{transform} CSS
(\lstinline|translate(...) scale(...)|). O zoom é \enquote{suave}: a roda do mouse
ajusta um \lstinline|targetScale| (fator \lstinline|exp(-deltaY*0.0016)|, clamp
0{,}1--5) e \lstinline|currentScale| é interpolado por \emph{requestAnimationFrame},
ancorado sob o cursor. Há também \emph{pan} por arraste, gestos de toque
(\emph{pinch} e \emph{pan}), \lstinline|fitToScreen()| (ajusta a $\sim$90\% do
\emph{viewport}), \lstinline|resetView()| e atalhos de teclado
(\lstinline|Ctrl+=/-/0/F|, \lstinline|Ctrl+R| recompila, \lstinline|F11| tela
cheia). \lstinline|downloadSVG()| serializa o SVG atual (clonado, com
\emph{namespaces} XML adicionados) para um \lstinline|Blob| baixável, nomeado pelo
módulo corrente --- funciona em qualquer nível de navegação.

\section{Simulação interativa com DigitalJS (O9)}
\label{sec:digitaljs}

Desde a fase O9, o \tool{PRISM} oferece, além do esquemático estático, uma
\emph{simulação lógica interativa} do mesmo \emph{netlist}, via \tool{DigitalJS}.
O botão \enquote{Simular} (id \lstinline|simToggle|) alterna entre as duas
superfícies. \textbf{A integração está de fato funcional} --- não é apenas uma
dependência presente: há \emph{handler} IPC dedicado, conversão real do
\emph{netlist} e renderização da paper interativa, conforme o código abaixo.

As dependências relevantes no \file{package.json} da \tool{AURORA} são
\lstinline|digitaljs| \lstinline|^0.14.2| e \lstinline|yosys2digitaljs|
\lstinline|^0.10.3| (instalado com \lstinline|main: dist/index.js| e um
sub-\emph{entry} \lstinline|core|).

\subsection{Construção do circuito no main}

O \emph{handler} \lstinline|prism:build-digitaljs| (em \file{main/ipc/prism.js})
resolve o top-level do \file{.spf} e chama
\lstinline|buildDigitalJSCircuit(paths, topLevelModule, tempDir)|. Esse fluxo é
deliberadamente \emph{separado} do esquemático para nunca perturbá-lo:

\begin{enumerate}[leftmargin=2.2em]
  \item \lstinline|collectSynthFiles| coleta os \file{.v} com a mesma lógica de
        \lstinline|runYosysCompilationWithPaths| (HDL + \lstinline|synthesizableFiles|
        + \path{TopLevel/} + \path{<proc>/Hardware/}).
  \item Gera um \emph{script} \tool{Yosys} \enquote{amigável ao
        yosys2digitaljs}, em nível de palavra (sem \lstinline|abc|/\lstinline|techmap|,
        para manter \lstinline|$add|/\lstinline|$mux|/\lstinline|$dff| como
        células \tool{DigitalJS} em vez de \emph{gates}):
\end{enumerate}

\begin{lstlisting}[caption={Script Yosys da simulação (buildDigitalJSCircuit)}]
read_verilog "..."
...
hierarchy -top <topLevelModule>
proc
opt_clean
memory -nomap
wreduce -memx
opt_clean
write_json "digitaljs.json"
\end{lstlisting}

\begin{enumerate}[leftmargin=2.2em,start=3]
  \item Roda o \emph{próprio} \tool{Yosys} (allowlisted) via
        \lstinline|spawnTracked|, com \emph{timeout} de 45\,s (uma síntese
        descontrolada falha com mensagem clara em vez de \emph{spinner} infinito).
  \item Conta as células do \emph{netlist}; se passar de \lstinline|MAX_CELLS = 3000|,
        \emph{recusa} a simulação com orientação para usar o esquemático
        (\tool{DigitalJS} é um simulador didático para módulos pequenos; o
        conversor roda síncrono no processo principal e o \emph{layout} no
        navegador degrada além de alguns milhares de células).
  \item Carrega \emph{lazy} o conversor \lstinline|require('yosys2digitaljs/core')|
        e chama \lstinline|yosys2digitaljs(yosysJson, {})| --- uma função
        \emph{pura} que recebe o JSON do \tool{Yosys} já produzido (nenhum
        \emph{spawn} adicional de \tool{Yosys}) e devolve o \lstinline|TopModule|
        no formato \tool{DigitalJS}.
  \item \lstinline|stripYosysLabels(circuit)| zera \emph{labels} internos
        gerados pelo \tool{Yosys} (que começam com \lstinline|$|, como
        \lstinline|$xor$src:line$n|), preservando os úteis (\lstinline|a|,
        \lstinline|b|, \lstinline|sum| etc.).
\end{enumerate}

O \emph{handler} retorna \lstinline|{ ok: true, circuit, topLevelModule }|, com
\emph{logs} de progresso por fase (\enquote{synthesizing\dots}, \enquote{netlist
N cells --- converting\dots}, \enquote{converted to M devices\dots}).

\subsection{Renderização interativa no renderer}

No renderer, \lstinline|enterSimMode()| (em \file{html/prism/prism.js}):

\begin{itemize}[leftmargin=1.4em]
  \item Usa um \emph{guard} de reentrância (\lstinline|_simBusy|) para que um
        segundo clique não inicie um segundo \emph{build} concorrente.
  \item Obtém o circuito via \lstinline|electronAPI.buildDigitalJS(paths)|.
  \item \emph{Carrega \tool{DigitalJS} de forma \emph{lazy}}
        (\lstinline|_loadDigitalJS()|): expõe o \lstinline|jQuery| global,
        importa o \lstinline|jquery-ui| \emph{completo} (na ordem correta, para
        que \lstinline|$.widget|/\lstinline|$.fn.dialog| existam antes de o
        \emph{bundle} do \tool{DigitalJS} avaliar seu \lstinline|dialog.js|), e
        só então faz \lstinline|import('digitaljs')| para obter a classe
        \lstinline|Circuit|. Manter isso fora do \emph{import} de topo significa
        que o fluxo do esquemático nunca avalia o \tool{DigitalJS} (um erro de
        carga não derruba a janela inteira) e o \emph{chunk} de $\sim$2\,MB só é
        baixado sob demanda.
  \item Instancia \lstinline|new Circuit(res.circuit, { layoutEngine: 'dagre' })|
        e exibe com \lstinline|circuit.displayOn(paperHost)|, depois
        \lstinline|circuit.start()|. O motor de navegador \emph{síncrono} (padrão)
        com \emph{layout} \lstinline|dagre| evita qualquer \emph{Web Worker},
        permitindo rodar sob \lstinline|file://|.
\end{itemize}

A simulação tem seu próprio \emph{pan}/\emph{zoom} (transform CSS no
\lstinline|.djs-wrapper|, espelhando os valores do esquemático), \emph{pan} por
arraste de área vazia (\lstinline|blank:pointerdown|), \lstinline|_fitPaper()|
para centralizar/ajustar, e \lstinline|_buildValueOverlays()|, que sobrepõe um
dígito vivo \lstinline|0|/\lstinline|1|/\lstinline|x| em cada caixa de
entrada/saída de 1 bit (o \tool{DigitalJS} sozinho só pinta a caixa) e o atualiza
a cada mudança de sinal. \lstinline|exitSimMode()| e \lstinline|_destroyCircuit()|
param e descartam o circuito e a \emph{paper} JointJS, evitando vazamentos em
ciclos repetidos de entrar/sair. Uma recompilação do esquemático sai
automaticamente do modo de simulação (o circuito vivo fica obsoleto).

\section{Árvore de hierarquia pós-síntese}

A \enquote{árvore de hierarquia} que o \tool{PRISM} apresenta não é uma estrutura
em árvore separada, e sim a \emph{navegação} sobre os módulos individuais que
\lstinline|splitHierarchyJson| derivou do \path{hierarchy.json}. Como o
\emph{script} do esquemático roda \lstinline|hierarchy -top| seguido de
\lstinline|proc| (sem \emph{flatten}), o \path{hierarchy.json} preserva a
estrutura hierárquica: cada módulo guarda suas células, e células cujo
\lstinline|type| é outro módulo clicável tornam-se pontos de entrada.

Em tempo de execução, o usuário começa no top-level e \enquote{desce} a árvore
clicando em células de submódulo; os \emph{breadcrumbs} mostram o caminho atual e
permitem saltar para qualquer ancestral. Como cada nível corresponde a um
\path{<modulo>.json} já materializado, descer/subir é instantâneo --- nenhuma
reexecução de \tool{Yosys} ocorre durante a navegação (apenas uma nova chamada a
\lstinline|netlistsvg.render| por módulo visitado).

\section{Resumo de artefatos e responsabilidades}

\begin{tabularx}{\linewidth}{Y Z}
\toprule
\textbf{Arquivo} & \textbf{Responsabilidade} \\
\midrule
\file{main/ipc/prism.js} & janela, \emph{pipeline} Yosys, split por módulo, render SVG, skins, DigitalJS; \emph{handlers} IPC \\
\file{js/app/preload\_prism.js} & \emph{bridge} sandboxizada com \emph{allowlist} de canais \\
\file{html/prism/prism.html} & layout da janela (toolbar, breadcrumbs, canvas SVG, host DigitalJS) \\
\file{html/prism/prism.js} & classe \lstinline|PRISMViewer|: i18n, SVG, navegação, pan/zoom, sim \\
\file{js/compilation/compilation\_flow.js} & botão PRISM, pré-flight, payload de caminhos, \emph{overrides} \\
\file{main/render\_loader.js} & seleção dev/prod do renderer (\lstinline|loadPage|) \\
\repo{netlistsvg-aurora} & fork de \lstinline|@silimate/netlistsvg| (fix de tsconfig; skins default) \\
\file{assets/prism-skins/*.svg} & símbolos customizados da família SAPHO (mesclados em runtime) \\
\file{scripts/prism-skin-standard.js} & gerador do padrão de símbolos SAPHO \\
\bottomrule
\end{tabularx}

\begin{nota}
Em síntese: o \tool{PRISM} sintetiza o RTL do projeto com o \tool{Yosys} bundle
(\lstinline|read_verilog -setattr src| $\rightarrow$ \lstinline|hierarchy| $\rightarrow$
\lstinline|proc| $\rightarrow$ \lstinline|setundef -zero| $\rightarrow$
\lstinline|opt_clean| $\rightarrow$ \lstinline|write_json|), divide o
\emph{netlist} por módulo, e renderiza cada módulo \emph{in-process} com o fork
\repo{netlistsvg-aurora} usando \emph{skins} SAPHO mescladas em tempo de
execução, produzindo um esquemático SVG navegável. Sobre o mesmo desenho de
fluxo, a fase O9 acrescentou uma simulação lógica interativa via
\tool{yosys2digitaljs}~+~\tool{DigitalJS}, integrada e funcional, ativada pelo
botão \enquote{Simular}.
\end{nota}
